\chapter{Conclusion}

\label{chap:conclusion}

This dissertation develops the reverse-engineering and measurement
infrastructure required to study pentimenti-class side-channels in
commercial ReRAM. The central claim has been that observable write
latency in commercial ReRAM decomposes into a structural component
controlled by polarity serialization and embedded ECC activity, and a
stochastic component controlled by conductive filament dynamics, and
that both components are recoverable through non-invasive timing
analysis alone. The three works compiled in this dissertation establish
that claim, and in doing so they expose a defect-dynamics measurement
channel in a device family for which no such channel was previously
available.

\section{Summary of Contributions}

Chapter~\ref{chap:fpga-sensor} closes a vendor-specificity gap in FPGA
Pentimenti measurement. By reformulating fine-grained sensor tuning as a
constrained PLL configuration search, the dual-edge TDC achieves an
average phase-step granularity of 1.6ps on Intel Cyclone~V, two orders
of magnitude finer than the vendor-default 78.125ps floor and sufficient
to resolve the picosecond-scale delay shifts characteristic of
BTI-induced pentimenti. This broadens the FPGA platforms on which the
attack class can be studied and removes the prior dependence on
Xilinx-specific phase-shift primitives.

Chapter~\ref{chap:brewrc} presents the first systematic characterization
of write latency in commercial ReRAM crossbar memories and introduces
BREW-RC, the first non-invasive technique for recovering the bit-exact
ECC parity submatrix of a commercial ReRAM product from write timing
alone. The four-experiment characterization identifies polarity
serialization and ECC parity activity as the dominant sources of
structured timing variance, motivates the directional Hamming distance
as the natural unit of analysis, and produces the labelled observation
classes from which BREW-RC's SAT formulation is derived. The recovered
parity submatrices improve Pearson correlation between predicted and
measured write latency from $0.203$ to $0.998$ on the MB85AS4MT and from
$0.180$ to $0.553$ on the MB85AS8MT, and are validated bit-for-bit
against an independent destructive reverse-engineering result on the
MB85AS4MT.

Chapter~\ref{chap:trng} applies the structural understanding from
Chapter~\ref{chap:brewrc} to entropy harvesting in commercial
off-the-shelf ReRAM. Subtracting the codeword-DHD timing model from
measured single-byte transition latencies eliminates the Sierpinski-like
structural pattern entirely, confirming that ECC activity is the sole
source of structured variance in these devices and that prior entropy
claims, which conflated ECC-driven perturbations with stochastic
behavior, are correspondingly overstated. ECC-aware construction of
wear-leveled, fixed-composition codeword sequences enables rigorous
min-entropy estimation under NIST SP~800-90B, and the resulting TRNG
passes SP~800-22 and PractRand while achieving a $17\times$ throughput
improvement over prior ECC-unaware methods.

Most importantly, the entropy-versus-wear characterization in
\S\ref{sec:trng-results} demonstrates that under fixed codeword DHD,
write-latency min-entropy varies systematically with per-bit cycling
history. Because DHD is held constant, the observed entropy drift cannot
be attributed to changes in switching composition, it reflects changes
in the conductive filament itself. This is concrete evidence that the
isolated stochastic component of write latency carries information about
cumulative defect dynamics, the property that pentimenti-class analysis
requires.

\section{Implications for Deployed ReRAM Products}

The findings of this dissertation have two immediate consequences for
the security of deployed commercial ReRAM.

First, the undocumented ECC structure of commercial ReRAM is not a
defense. BREW-RC recovers the bit-exact parity submatrix from black-box
write-access timing in a matter of hours of measurement and a parallel
SAT search bounded by $k$. Any security argument that implicitly relied
on the stochasticity of commercial ReRAM without ECC-awareness,
including entropy claims for prior ReRAM TRNGs that did not account for
codeword-level switching, must be revisited.

Second, the write latency exposed at the memory interface leaks useful
information. The DHD of the underlying codeword, not the user-visible
data word, determines latency distribution membership. An adversary with
knowledge of the ECC can recover information about the codeword being
programmed, and with sufficient measurement budget can characterize the
stochastic component associated with individual bitcells. Whether this
is sufficient to recover data remanence from a prior user (the ReRAM
analog of the FPGA Pentimenti attack) remains an open question.

\section{Toward ReRAM Pentimenti}
This dissertation establishes the measurement infrastructure for
pentimenti analysis in commercial ReRAM, but it does not demonstrate
ReRAM Pentimenti recovery itself. The wear-induced entropy drift result
is necessary evidence that filament defect dynamics leak through the
isolated stochastic timing component, but does not yet establish a
recoverable signal analogous to the BTI-based component exploited by
FPGA Pentimenti. A ReRAM pentimenti attack requires a recoverable signal
that is data-dependent on the prior user's activity, persists, and
survives reads or reprogramming by the adversary. Each of these
properties must be characterized experimentally before the attack class
can be said to extend to commercial ReRAM. Measurement of filament
degradation is a first step toward this objective. Other measurement
defect dynamics, such as the relaxation dynamics described in
\S\ref{sec:bg-defects} suggest that the post-RESET HRS distribution
carries a measurable dependence on the cell's most recent SET history,
and that dependence is the natural ReRAM analog of the recoverable BTI
component exploited in FPGA Pentimenti, future work is required to
determine how this and similar pentimenti mechanisms could be detected
under the measurement framework presented in this work.
